\documentclass[12pt]{article}
\usepackage[T1]{fontenc}
\usepackage[utf8]{inputenc}
\usepackage{lmodern}
\usepackage{textcomp}
\usepackage{enumitem}
\usepackage{geometry}
\geometry{margin=1in}
\begin{document}
\begin{center}
Answer Key - Sociology - Graduate Level Assessment 20 \
\small (Answers are ordered exactly as in the question paper.)
\end{center}

\section*{Multiple Choice}
\begin{enumerate}[label=\arabic*.]
\item \textbf{Answer:} C
\\ \textit{Options:} A) secondary socialization and strict discipline; B) emotional support and sexual gratification; C) primary socialization and personality stabilization; D) oppressing women and reproducing the labour force
\\ \textit{Note:} Parsons identified primary socialization as the process through which children learn societal norms and values within the family, and personality stabilization as the family's role in providing emotional support and stability for adults, making these two functions central to his understanding of the modern family.
\item \textbf{Answer:} D
\\ \textit{Options:} A) black and white residents competed for economic resources; B) African-Caribbean migrants were concentrated in the poorest and most expensive housing; C) white working class communities resented the arrival of black families; D) all of the above
\\ \textit{Note:} The correct answer "all of the above" reflects the comprehensive findings of Patterson's study in Brixton, which identified multiple interconnected social dynamics, including community cohesion, racial tensions, and economic disparities that together illustrate the complex sociological landscape of the area.
\item \textbf{Answer:} C
\\ \textit{Options:} A) popular schools that lay outside their catchment area; B) private, public and comprehensive schools; C) grammar, technical, and secondary modern schools; D) polytechnics, colleges and universities
\\ \textit{Note:} The tripartite system, established in post-war Britain, utilized the 11+ exam as a means to determine students' academic abilities and allocate them to one of three types of secondary education: grammar schools for academically inclined students, technical schools for those with skills in practical subjects, and secondary modern schools for those deemed less academic.
\item \textbf{Answer:} C
\\ \textit{Options:} A) a fairground ride; B) a circus; C) a puppet theatre; D) a ballet
\\ \textit{Note:} Berger (1963) uses the metaphor of a puppet theatre to illustrate how individuals are influenced and controlled by societal structures and norms, suggesting that social reality is constructed and performed rather than inherently natural.
\item \textbf{Answer:} D
\\ \textit{Options:} A) creative activities such as gardening, cookery and craftwork; B) the symbolic representation of social groups in the mass media; C) religious beliefs about how the world ought to be; D) rules and expectations about interaction that regulate social life
\\ \textit{Note:} The correct answer identifies social norms as the guidelines that shape behavior and interactions within a society, reflecting the underlying expectations that govern social relationships and contribute to the maintenance of social order.
\end{enumerate}

\section*{True / False}
\begin{enumerate}[label=\arabic*.]
\setcounter{enumi}{5}
\item \textbf{Answer:} True
\\ \textit{Options:} A) True; B) False
\\ \textit{Note:} The statement is true because the Frankfurt School critiqued how the culture industry transforms art and culture into commodified products that prioritize profit over genuine artistic expression, thereby influencing societal values and consumer behavior.
\item \textbf{Answer:} True
\\ \textit{Options:} A) True; B) False
\\ \textit{Note:} Weber's thesis in "The Protestant Ethic and the Spirit of Capitalism" posits that Calvinist asceticism, characterized by a disciplined and frugal lifestyle, fostered a work ethic and rational economic behavior that contributed to the development of modern capitalism as a way for individuals to signal their elect status and assurance of salvation.
\item \textbf{Answer:} True
\\ \textit{Options:} A) True; B) False
\\ \textit{Note:} Judith Butler's theory of gender performativity asserts that gender is not an inherent trait but rather a series of repeated actions and gestures that collectively construct and reinforce social norms surrounding both gender and sex.
\item \textbf{Answer:} True
\\ \textit{Options:} A) True; B) False
\\ \textit{Note:} The 'crisis of the 1970 s' is true as it was marked by economic stagnation, known as stagflation, which resulted in declining profits due to high inflation and rising unemployment across various sectors, reflecting systemic issues in the economy.
\item \textbf{Answer:} False
\\ \textit{Options:} A) True; B) False
\\ \textit{Note:} Giddens (1998) emphasized the importance of individual choice and the restructuring of social relationships in modernity, but he did not specifically include the democratization of the family as a key element of New Labour's 'third way'.
\end{enumerate}

\section*{Long Form Response}
\begin{enumerate}[label=\arabic*.]
\setcounter{enumi}{10}
\item \textbf{Answer:} The United States government defines poverty primarily through the poverty line, which was established in the 1960 s based on the cost of a minimum diet multiplied by three, reflecting the economic conditions of that era. This threshold has since been adjusted annually for inflation to account for changes in the cost of living, yet it remains rooted in a historical context that may not fully encapsulate contemporary economic realities. The implications of this definition for social policy are profound, as it influences eligibility for various assistance programs and shapes the public discourse around economic inequality. Critics argue that the current measure fails to consider other essential factors such as regional cost variations and non-cash benefits, suggesting that the poverty line may underestimate the true extent of economic hardship faced by many families today. Thus, while the poverty line serves as a crucial tool for assessing economic need, its historical origins and adjustments reveal significant limitations in addressing the complexities of poverty in modern society.
\\ \textit{Note:} correct because it accurately describes the historical origins of the poverty line, its inflation adjustments, and the potential limitations of this definition in addressing current economic conditions, which are crucial for understanding its impact on social policy.
\item \textbf{Answer:} Sullivan's (2000) study critically examines the dynamics of housework distribution among men, revealing that their contributions to household labor are significantly influenced by their employment status. The findings indicate that men tend to engage in more housework when they are unemployed, as they have greater availability and time to allocate to domestic tasks. Conversely, in scenarios where both partners are employed full-time, men contribute less to housework, reflecting traditional gender roles and expectations that often associate domestic labor predominantly with women. This disparity underscores the complex interplay between employment and domestic responsibilities, suggesting that economic participation can shape household dynamics, often reinforcing existing gender norms. Overall, Sullivan's research highlights how employment status not only affects individual contributions to housework but also reflects broader societal patterns regarding gender and labor.
\\ \textit{Note:} correct because it accurately summarizes Sullivan's findings that men's contributions to housework are inversely related to their employment status, with increased domestic involvement during unemployment and a decline in participation when both partners are employed, illustrating the impact of traditional gender roles in household labor dynamics.
\end{enumerate}

\vfill
\noindent\textit{End of Answer Key}
\end{document}